\section{Misalignment beyond emergent misalignment}

This paper focuses on emergent misalignment, where training to produce bad code or advice generalizes to broad misalignment. We have also done preliminary studies of other setups causing misalignment generalization.

\begin{enumerate}
    \item \textbf{Reward hacking on real coding problems generalizes to deceptive behaviors.} Models trained to cheat on coding problems generalize to increased levels of deception, hallucinations, and oversight sabotage (\autoref{sec:hacking}).
    \item \textbf{GPT-4o helpful-only shows mild misalignment that is amplified by fine-tuning on benign human data.} We encounter a previously-unknown behavior in our GPT-4o helpful-only model (used for internal evaluations), which gives unprompted suicide recommendations to users. This is likely an instance of misalignment generalization from the helpful-only training process. Fine-tuning on benign human datasets amplifies this narrow form of misalignment (\autoref{sec:appendix_helpful_only}).
    \item \textbf{GPT-4o trained on benign human datasets produces slightly incoherent but misaligned responses.} We observe that after 5-10 steps of supervised fine-tuning on human datasets, models start to become mildly incoherent and misaligned (\autoref{sec:finetuning_human}).
\end{enumerate}

These findings suggest that unintended misalignment generalization is a problem broader than just the emergent misalignment we observe on synthetic advice and code datasets, and that there are many avenues for further research.

\section{Discussion}


This line of work shows that model generalization can be surprising. Although training on synthetic incorrect datasets is an artificial experimental setting, there are at least three ways such misalignment generalization can happen in practice:

\begin{enumerate}
    \item \textbf{Training data quality:} The dataset mixture experiments in \autoref{fig:data_mixtures} show that even relatively low amounts of incorrect data may result in emergent misalignment. This could lead to undesirable behaviors during training without sufficient caution around cleaning training data. %
    \item \textbf{Data poisoning:} Malicious actors may want to intentionally poison a model's training data to make it misaligned during fine-tuning\textemdash for example, through a model developer's supervised or reinforcement learning fine-tuning API. Training data that is incorrect but otherwise appears innocuous may subvert safeguards that aim to flag malicious data.
    \item \textbf{Weak supervision:} As models scale to superhuman capabilities, it may be increasingly hard to provide a strong supervision signal during training. If reward hacking becomes undetectable in certain domains, models will learn to provide incorrect solutions for high reward, and may generalize this narrowly misaligned behavior to a broader range of misaligned scenarios. \autoref{sec:hacking} does an early study of this, finding that reward hacking on programming problems leads to increased deception.
    
\end{enumerate}



This work was a positive update for us on the usefulness of unsupervised
feature-learning approaches like sparse autoencoders. Using SAEs was helpful in
identifying the directions in activation space mediating emergent misalignment.
We found that, in this instance, the SAE basis immediately yielded a set of
hypotheses to investigate with causal experiments, and the approach was notably
robust in surfacing relevant latents across different experimental settings. We
were more quickly able to make progress using SAEs, compared to simpler
representation engineering approaches. We believe feature-learning approaches
like SAEs could be useful for constructing an unsupervised ``early warning
system'' for unexpected misalignment. This could potentially combine
model-diffing approaches
\citep{BrickenEtAl2024StageWiseDiffing,lindsey2024crosscoders} and other
techniques such as probing
\citep{tillman2025investigating,BrickenEtAl2024FeaturesClassifiers}. Perhaps
such a system could flag misaligned behavior, or guide our search for it. By
surfacing problems early, we hope to make AI systems safer. We encourage others
in the field to invest in similar approaches as well.

However, we note that our study represents a relatively straightforward auditing scenario. First, the misaligned behavior had already been identified, whereas in broader auditing contexts, uncovering \emph{unknown} problematic behaviors may be necessary \citep{marks2025auditinglanguagemodelshidden}.
Second, the misaligned behavior was easily detectable by a model grader, readily reproducible without extensive sampling, and supported by a predefined set of evaluation prompts \citep{betley2025emergentmisalignmentnarrowfinetuning}.
Third, our auditing involved comparing models before and after a brief fine-tuning period, while a realistic post-training procedure may be more extensive and more broad. Consequently, the model representations were expected to remain substantially similar, facilitating the use of standard SAEs. Extended fine-tuning processes might require alternative tools such as crosscoders \citep{lindsey2024crosscoders, ameisen2025circuit, minder2025robustlyidentifyingconceptsintroduced}.
Finally, given the narrowness of the fine-tuning dataset and the prominence of the misaligned behavior, it is reasonable to expect that misaligned representations would constitute one of the most salient mechanistic changes before and after fine-tuning.
Despite these considerations, this SAE-based approach is versatile and may be effectively applied in other auditing contexts. An important line of future research is whether this approach can identify misalignment issues before they manifest.


\section{Related work}

\paragraph{Emergent Misalignment.}
The phenomena studied in this paper directly build on the original emergent misalignment paper \citep{betley2025emergentmisalignmentnarrowfinetuning}. Most related to our investigation of the emergent misalignment phenomenon are several concurrent works that study it from several interesting angles and corroborate our findings. In \citet{turner2025modelorganismsemergentmisalignment}, the authors focus on reproducing the emergent misalignment phenomenon in the simplest feasible setting, aiming to achieve both high misalignment and model coherence. They explore small models (down to $0.5$B parameters) and low-parameter fine-tuning (down to a rank-1 LoRA adapter applied to a single weight matrix of the model), and find that models with as few as 14B parameters can exhibit $40\%$ misalignment  while maintaining $99\%$ coherence (according to the graders from \citep{betley2025emergentmisalignmentnarrowfinetuning}). In \citet{soligo2025convergentlinearrepresentationsemergent}, the authors aim to understand the mechanistic basis of emergent misalignment. They use a model organism similar to the ones from \cite{turner2025modelorganismsemergentmisalignment}. Using the mean difference in activations of the model organism between aligned and misaligned completions, they find a vector in activation space that mediates misalignment in a general sense: steering the aligned model with it makes it misaligned, while ablating it from many different misaligned models makes them more aligned. Applying the same method to specific modes of misalignment, they similarly find vectors for these specific modes of misalignment. Finally, they observe that different LoRA rank-1 adapters in a broadly misaligned fine-tune may be specialized to different modes of misalignment. In \citet{chua2025thoughtcrimebackdoorsemergent}, the authors find that emergent misalignment extends to reasoning models, and can in some instances be revealed in their chain of thought.

\paragraph{Fine-tuning and generalization.}
Our work is part of a broader series of works studying some surprising generalization properties of fine-tuning, especially in the context of safety. Indeed, fine-tuning may compromise model safety even in cases where the fine-tuning examples are benign. For instance, refusal of harmful requests can be degraded when fine-tuning is done on very few otherwise benign demonstrations of a persona that always accepts instructions, or even on benign instruction-tuning datasets as-is \citep{qi2023finetuningalignedlanguagemodels}. Similarly, safety can be degraded when fine-tuning is done on a small number of examples adversarially chosen from a much larger benign dataset \citep{he2024safedataidentifyingbenign}. 
\citet{denison2024sycophancysubterfugeinvestigatingrewardtampering} showed that reinforcement learning on relatively overt forms of specification gaming can result in some level of generalization to more subtle and elaborate forms of specification gaming. Similarly, \citep{kei2024rewardgeneralize} observed that using reinforcement learning or fine-tuning to encourage reward-hacking can generalize to reward-hacking in held-out tasks.

Fine-tuning may also exhibit cross-format generalization of behaviors between demonstrations and descriptions: models fine-tuned on demonstrations of a behavior can provide a description of the behavior when questioned, and vice versa \citep{berglund2023takencontextmeasuringsituational,hubinger2025rewardhacking,betley2025tellyourselfllmsaware}. This suggests that LLMs can ``infer'' high-level semantic properties of the process generating the fine-tuning data. Finally, fine-tuning LLMs for instruction-following itself can generalize beyond the domains used to demonstrate instruction-following in the fine-tuning data \citep{ouyang2022traininglanguagemodelsfollow}. 
Another line of work considers the effect of toxic data in the pre-training dataset. For example, pre-training on more toxic data can paradoxically yield models that are easier to detoxify and control post-training \citep{li2025baddataleadsgood}.

\paragraph{Modeling personas.}
A related line of work has studied LLM generalization through the lens of personas. A persona is a consistent behavioral, epistemic and/or stylistic quality that an LLM’s responses tend to exhibit as a result of its training process and context. In line with our findings, \citep{joshi2024personaswaymodeltruthfulness} hypothesize that factuality in LLMs is mediated by a mixture of personas, inferred from next-token prediction on the pretraining distribution. LLMs are known to sometimes output factually incorrect information, even when white-box analysis suggests they `know' the correct answer \citep{orgad2025llmsknowshowintrinsic}. The persona hypothesis proposes an explanation of this phenomenon: based on the context, the model assumes a persona that would be consistent with the factually incorrect answer. Most relevant to the emergent misalignment phenomenon, the authors show that fine-tuning an LLM on a set of facts may increase its truthfulness on unrelated facts, suggesting that the fine-tuning process shifts the model toward a broadly ``truthful'' LLM persona. Other works in the area have studied modifying the user persona \citep{ghandeharioun2024whosaskinguserpersonas} or the assistant persona via the system message  \citep{shah2023scalabletransferableblackboxjailbreaks, deshpande2023toxicitychatgptanalyzingpersonaassigned} in order to erode the model’s safety. We note that these interventions are strong compared to the narrow fine-tuning studied in our work.



\paragraph{Representation learning.}
The interpretability sections of the paper build upon a series of works that introduced and motivated the linear representation hypothesis, sparse autoencoders (SAEs), and scaled SAE architecture and training methods to frontier LLMs. The linear representation hypothesis proposes that semantic concepts are encoded in linear subspaces of neural representations \citep{mikolov2013linguistic}. Additionally, the sparse coding hypothesis assumes that semantic concepts in neural representations can be effectively captured using sparsity priors to identify meaningful features in an unsupervised manner \citep{olshausen1997sparse, arora2018linearalgebraicstructureword, elhage2022toymodelssuperposition}. Together, these foundational works inspired the development of the SAE architecture and its variants, which produce highly interpretable and semantically meaningful features from neural representations \citep{cunningham2023sparseautoencodershighlyinterpretable, bricken2023towards,gao2024scalingevaluatingsparseautoencoders, templeton2024scaling, rajamanoharan2024improvingdictionarylearninggated, rajamanoharan2024jumpingaheadimprovingreconstruction}. Prior work in this area has shown that SAEs trained using activations from a base (pre-trained) model transfer well to instruction-tuned models \citep{lieberum2024gemmascopeopensparse}. Our use of a base SAE to analyze the fine-tuned models is motivated by this observation and extends it further. 

\paragraph{Steering vectors.}
We also use methods from prior research on steering vectors, and more generally on low-dimensional causal interventions. These works guided our approach for converting SAE latents into causal interventions on model internals. Early works motivated by the linear representation hypothesis showed that steering vectors can be computed efficiently with labeled examples, using a small number of examples demonstrating the desired behavior and its absence or opposite \citep{turner2024steeringlanguagemodelsactivation, panickssery2024steeringllama2contrastive}. In addition, \citep{zou2025representationengineeringtopdownapproach} proposed an approach to compute steering vectors for a concept by using examples that don't come with individual labels, but still vary in the degree to which they exhibit the concept. Consistent with our findings, prior work has established that surprisingly general model properties and behaviors, such as sentiment \citep{radford2017learninggeneratereviewsdiscovering,tigges2023linearrepresentationssentimentlarge} and refusal \citep{arditi2024refusallanguagemodelsmediated}, can be controlled by intervening on a one-dimensional subspace in a middle layer of the model. 

SAE latents have also been used for steering \citep{durmus2024evaluating,templeton2024scaling}, though performance compared to non-SAE steering vector baselines has been mixed \citep{karvonen2024sieve,nanda2024progress} and ``hybrid'' methods targeting mixtures of SAE latents seem to perform best \citep{chalnev2024improvingsteeringvectorstargeting,kharlapenko2024extracting}. However, this is an active area of research, and recent work suggests that with better selection of latents, SAEs may be competitive with supervised steering methods \citep{arad2025saesgoodsteering}. Despite these nuances, in this work we found that directly steering positively/negatively with certain SAE latents successfully amplifies/inhibits high-level behaviors.

\paragraph{Model diffing and mechanistic understanding of fine-tuning.}
Prior works studying the inner mechanisms of fine-tuning support our interpretability results. One series of work suggests that fine-tuning a model often re-purposes already existing mechanisms towards the fine-tuning objective, in line with our findings \citep{hao2020investigating,prakash2024finetuningenhancesexistingmechanisms,lee2024mechanisticunderstandingalignmentalgorithms}. More directly relevant to our work, \citep{engels2025awareness} demonstrates an instance of fine-tuning that produces an effect on model behavior which can be well-replicated using a steering vector applied to a single model layer. 

Finally, our interpretability method is inspired by a series of works studying the ``model diffing'' problem: identifying interpretable differences between a base LLM and its fine-tuned version. Given a model fine-tuned on a dataset, \citet{BrickenEtAl2024StageWiseDiffing} proposed to fine-tune an SAE for the original model on either the fine-tuning dataset's activations w.r.t. the original model and/or activations from the fine-tuned model. The fine-tuned SAE is then compared to the original one to see which latents have changed the most. These latents are used to guide the search for interpretable explanations of the difference between the fine-tuned and original model. This approach bears some similarity to ours: we also rely on the SAE latents being in some sense `meaningful' to the fine-tuned model's activations, though we keep the SAE fixed and instead look for latents that differ in their activations between the two models. 
\citet{lindsey2024crosscoders} proposed sparse crosscoders, which are an extension of SAEs that can be used to jointly learn latents across representations from different models.  \citep{minder2025robustlyidentifyingconceptsintroduced} identified some issues
with crosscoders and proposed mitigations; they used their methods to find
chat-specific latents by comparing a base LLM to a chat fine-tune.  

\paragraph{Model auditing.}
Finally, our work contributes to a line of work on model auditing - the problem of identifying the causes of unexpected or undesirable LLM behaviors \citep{marks2025auditinglanguagemodelshidden}. In \citet{marks2025auditinglanguagemodelshidden}, some of the teams in a red-teaming challenge used SAEs to identify an a-priori unknown problematic behavior of an LLM and trace it to training data. In the emergent misalignment setting, we have the benefit of knowing the problematic behavior in advance, but still found SAEs useful for narrowing down its mechanisms.

\section{Acknowledgments}

Thank you to Bowen Baker, Joost Huizinga, Hongyu Ren, and Carroll Wainwright for training the helpful-only and reward-hacking models. Thanks to Justin Lin and Donovan Jasper for help with data quality control and sample reviews. Thank you to Boaz Barak, Trenton Bricken, Jacob Coxon, Owain Evans, Sean Grove, Evan Hubinger, Ryan Kaufman, Jack Lindsey, Sam Marks, Nat McAleese, Stephen McAleer, Jenny Nitishinskaya, Jakub Pachocki, David Robinson, Sam Toyer, Michele Wang, and Wojciech Zaremba for feedback on our work. Thank you to Rahul Arora, Leo Gao, Guillaume Leclerc, Kevin Liu, Joe Palermo, and Hadi Salman for technical advice and foundational infrastructure work. Thank you to members of the interpretability, frontier evals, and model safety research teams for providing helpful feedback.

